\subsection{3. Proposed A2C-Based
	Framework}\label{proposed-a2c-based-framework}}

\hypertarget{problem-formalization}{%
	\subsubsection{3.1 Problem Formalization}\label{problem-formalization}}

We formalize the problem following the definitions from prior work
{[}1{]}:

\textbf{Definition 1: Moving Crowdsourced Service MS} \cite{neiat2021}. A moving
crowdsourced service MS is a tuple of \(\langle id, F, Q \rangle\)
where: - \(id\) is a unique service identifier - \(F\) is a function
offered by MS (e.g., providing a moving WiFi hotspot). The function
represents a moving service's coverage in space and time, defined as a
moving region \(\langle T_s, R_{ti}(p_i) \rangle\) where: -
\(T_s = \{\langle t_i, x_i, y_i \rangle\}\) is a service trajectory---a
sequence of timestamped samples where \((x_i, y_i)\) is
longitude/latitude at timestamp \(t_i\) - \(R_{ti}(p_i)\) is the
coverage region offered by MS at time \(t_i\), represented as a circular
area centered at \(p_i\) with radius \(r\) - \(Q\) is a set of QoS
attributes (e.g., capacity)

\textbf{Definition 2: User Trajectory \(T_u\)} \cite{neiat2021}. A user trajectory is the
path traveled by a user, defined as a set of \(k\) timestamped samples:
\[T_u = \{\langle u_{t_i}, u_{x_i}, u_{y_i} \rangle\}\]

\textbf{Definition 3: Spatial Candidate Pair} \cite{neiat2021}. Given a set of moving
services \(\mathcal{M} = \{MS_1, MS_2, ..., MS_n\}\), a user trajectory
\(T_u = \{up_1, up_2, ..., up_n\}\) where \(up_i = (u_{x_i}, u_{y_i})\),
and a search radius \(r_s\), a moving service forms a spatial candidate
pair \(cp_t^i\) for timestep \(t_i\) if its location \(MS_k.p_i\) at
\(t_i\) is inside a disk region \(D_{t_i}\) (center = \(up_i(t_i)\),
radius = \(r_s\)):
\[\text{Valid}_{spatial}(i, t_i) = \mathbb{1}(d(T_u.up_{t_i}, MS_k.p_{t_i}) \leq r_s)\]
where \(d(\cdot)\) is the Euclidean or Haversine distance.

\textbf{Definition 4: Valid Candidate Moving Service} \cite{neiat2021}. A moving service
\(MS_i\) is a valid candidate service for a given user trajectory if it
is paired with the user trajectory over \(w\) consecutive timesteps
where \(w > 0\):
\[CMS_i = \{cp_{t_a}, cp_{t_b}, ..., cp_{t_w}\}, t_a < t_b < ... < t_w, |a - b| = 1\]

\textbf{Problem Definition} \cite{neiat2021}: Given a set of moving crowdsourced services
\(\mathcal{M} = \{MS_1, MS_2, ..., MS_n\}\), a user trajectory \(T_u\),
and a search radius \(r_s\) as input, the problem is to find the
``optimal'' composition plan that gives the best trade-offs among
multiple QoS criteria---high QoS while maintaining a low number of
disconnections. The output is a composition plan \(CP\) which is a
sequence of moving crowdsourced services:

\begin{equation*}
CP = \{S_1, S_2, \ldots, S_n\} \quad \text{where} \quad S_i \in \mathcal{V}_u
\end{equation*}

\noindent where \(\mathcal{V}_u\) denotes the set of valid candidate
services for user trajectory \(T_u\).

\textbf{Assumptions} \cite{neiat2021}: - One moving service can only serve one user at
any point in time - A moving service moves between any two consecutive
timestamps \(t_i\) and \(t_{i+1}\) with constant speed, enabling
position interpolation in interval \([t_i, t_{i+1}]\) - Radii of all
coverage regions are fixed to a single value - Maximum spatial proximity
equals the radius of fixed WiFi hotspot coverage (e.g., 20-100m) - We
focus on deterministic moving crowdsourced services

\textbf{MDP Formulation}: We formulate this as an MDP, following the exact structure from prior work [1].

\textbf{State Space ($S$)}: At time step $t$, the state consists of the current user trajectory sample and all available service trajectories. Each sample is a tuple: $<t, x, y>$ for user trajectories, and $<t, x, y, service\_id, QoS>$ for service trajectories. At any given time, the environment has a current state $s \in S$ representing the current user trajectory sample reported to the agent.

\textbf{Action Space (\(A\))}: Select service ID from available
services: \[a_t \in \{1, 2, 3, ..., n\}\]

\textbf{Spatio-Temporal Validity}: A service is valid (spatial candidate
pair) at time \(t\) if:
\[\text{Valid}(i, t) = \mathbb{1}(d_{ij}(t) \leq R_{comm})\]

where \(d_{ij}(t)\) is the Euclidean distance between service \(i\) and
consumer \(j\) at time \(t\), and \(R_{comm}\) is the communication
range (discovery zone).

The capacity \(C_{ij}(t)\) derives from Euclidean distance through the
STR model. The agent learns to select services that are both within
communication range AND provide optimal capacity.

\textbf{Transition Dynamics (\(P\))}: Transitions follow the stochastic
dynamics of moving services and device mobility. Service positions
evolve according to their movement patterns observed in the trajectory
data. Connectivity depends on spatial proximity within communication
range \(R_{comm}\). The distance matrix \(D_t\) updates accordingly
with each transition.

\textbf{Reward Function (\(R\))}: The reward function is based on the
capacity QoS parameter. Since QoS attributes (capacity) are computed
from the distance between the consumer and moving service---which is
unknown a priori---the algorithm must learn the optimal execution
policy through interaction with the environment.

The reward for selecting service \(i\) at time \(t\) is:
\[r(s_t, a_t) = C_{ij}(t)\]

where \(C_{ij}(t) = B \cdot \log_2(1 + STR(d_{ij}(t)) \cdot SNR_{max})\)
is the STR-derived capacity, and \(d_{ij}(t)\) is the Euclidean distance
at time \(t\).

If the selected service is outside the discovery zone
(\(d_{ij}(t) > R_{comm}\)), a penalty is applied:
\[r(s_t, a_t) = \begin{cases} C_{ij}(t) & \text{if } d_{ij}(t) \leq R_{comm} \\ -1 & \text{if } d_{ij}(t) > R_{comm} \end{cases}\]

\textbf{Training Mode}: The experiments use \textbf{offline training}
mode where the agent learns from pre-collected trajectory data without
active environment interaction. In offline mode, the environment
provides fixed state-action-reward tuples from the dataset, enabling
sample-efficient learning from historical trajectories.

\textbf{Training/Test Split}: We use 70\% of the data in each dataset
for training and the remaining 30\% for testing.

\hypertarget{environment-interaction-and-training-algorithm}{%
	\subsubsection{3.3 Environment Interaction and Training
		Algorithm}\label{environment-interaction-and-training-algorithm}}

The training process follows the interaction paradigm between agent and
environment:

\textbf{Interaction Loop}: 1. Agent requests initial state from
environment (step a) 2. Environment fetches the current user trajectory
sample and reports it as current state (step b.1) 3. Environment fetches
the list of service IDs and reports them as possible actions (step b.2)
4. During exploration, agent randomly invokes an action by selecting a
valid service ID (step c) 5. Environment computes reward based on the
invoked action (step d.1) 6. Environment updates its current state to
the next user trajectory sample (step d.2) 7. Next state and reward are
sent to agent (step e) 8. Agent continues until all samples in user
trajectory are visited 9. Upon traversing all samples, environment
resets to first sample, process repeats

\textbf{Handling Edge Cases} \cite{neiat2021}:

\emph{Case 1 - No valid candidate services}: When no moving services
overlap with the current user trajectory sample in space and time, we
introduce a \textbf{dummy service}. Selecting the dummy service yields a
lower reward (-1) compared to normal rewards {[}0-1{]}, diverting the
agent from selecting invalid candidates.

\emph{Case 2 - Agent selects invalid service}: When the agent selects a
moving service that does not overlap with the current user trajectory
sample (either in time or space), we apply a penalty (-10). This ensures
the agent always favors the dummy service over invalid selections since
-1 \textgreater{} -10.

\begin{table}[htbp]
\centering
\caption{A2C Training Algorithm for Moving IoT Service Composition}
\label{a2c-training-algo}
\resizebox{0.9\linewidth}{!}{
\begin{tabular}{|l|p{10cm}|}
\hline
\textbf{Step} & \textbf{Operation} \\
\hline
1 & Initialize policy parameters $\theta_\\pi$ and value parameters $\theta_V$ \\
\hline
2 & Initialize replay memory $\\mathcal{D}$ with capacity $N$ \\
\hline
3 & Set exploration rate $\\epsilon \gets 1.0$ \\
\hline
4 & \textbf{for} episode $= 1$ \textbf{to} $M$ \\
\hline
5 & \quad Sample user trajectory $T$ from training set \\
\hline
6 & \quad \textbf{for} each timestep $t$ \textbf{in} $T$ \\
\hline
7 & \quad\quad Observe state $s_t$ from environment \\
\hline
8 & \quad\quad \textbf{if} rand() $< \epsilon$ \textbf{then} select random $a_t$ \textbf{else} $a_t \gets \\arg\\max_a \\pi(a|s_t;\\theta_\\pi)$ \\
\hline
9 & \quad\quad Execute $a_t$, observe reward $r_t$ and next state $s_{t+1}$ \\
\hline
10 & \quad\quad Store transition $(s_t, a_t, r_t, s_{t+1})$ in $\\mathcal{D}$ \\
\hline
11 & \quad\quad Sample mini-batch $\\mathcal{B}$ from $\\mathcal{D}$ \\
\hline
12 & \quad\quad Update critic: minimize TD-error on $\mathcal{B}$ \\
\hline
13 & \quad\quad Update actor: maximize advantage on $\mathcal{B}$ \\
\hline
14 & \quad \textbf{end for} \\
\hline
15 & \quad $\epsilon \gets \epsilon \times \gamma_\epsilon$ (decay exploration) \\
\hline
16 & \textbf{end for} \\
\hline
\end{tabular}}
\end{table}

\textbf{Edge Server Assumptions}: Model training and storage are carried
out using edge servers. Edge servers are assumed to be conveniently
accessible by moving IoT services. Each edge server serves a small
subset of moving devices, making storage and processing overheads
negligible.

\hypertarget{str-based-selection-with-capacity-integration}{%
	\subsubsection{3.2 STR-Based Selection with Capacity
		Integration}\label{str-based-selection-with-capacity-integration}}

The core selection function from prior work integrates into A2C using
the hierarchical STR → Capacity → Reward model:

\textbf{Step 1 - Euclidean Distance Calculation}: For each service \(i\)
and user/device position \(j\):
\[d_{ij} = \sqrt{(x_i - x_j)^2 + (y_i - y_j)^2}\]

\textbf{Step 2 - STR Calculation}:
\[STR(d_{ij}) = \begin{cases} 1 & \text{if } d_{ij} \leq R_c \\ e^{-k \cdot (d_{ij} - R_c)} & \text{if } d_{ij} > R_c \end{cases}\]

\textbf{Step 3 - Capacity Calculation}:
\[C_{ij} = B \cdot \log_2(1 + STR(d_{ij}) \cdot SNR_{max})\]

\textbf{Step 4 - Service Ranking}: Each service is ranked based on:
\[Score_{selection}(i) = C_{ij} \cdot E_i \cdot T_{available}^i\]

This hierarchical model ensures the selection function follows: distance
→ STR → capacity → reward, preserving the exact same approach from Paper
17.

\hypertarget{a2c-algorithm-design}{%
	\subsubsection{3.3 A2C Algorithm Design}\label{a2c-algorithm-design}}

The A2C algorithm maintains policy \(\pi(a_t|s_t; \theta_\pi)\) and
value function \(V(s_t; \theta_V)\) approximated by neural networks. The
advantage function guides policy updates:

\[A(s_t, a_t) = r_t + \gamma V(s_{t+1}; \theta_V) - V(s_t; \theta_V)\]

The policy gradient updates the actor parameters:
\[\nabla_\theta J = \mathbb{E}[A(s_t, a_t) \nabla_\theta \log \pi(a_t|s_t; \theta_\pi)]\]

The value function minimizes:
\[L_V = \mathbb{E}[(r_t + \gamma V(s_{t+1}) - V(s_t))^2]\]

The combined loss includes policy and value components with entropy
regularization: \[L_{total} = L_V + c_1 L_\pi - c_2 H(\pi)\]

where \(H(\pi)\) is the entropy of the policy distribution, encouraging
exploration.

\hypertarget{convergence-and-complexity-analysis}{%
	\subsubsection{3.3.1 Convergence and Complexity
		Analysis}\label{convergence-and-complexity-analysis}}

We analyze theoretical properties based on recent advances in
actor-critic convergence theory {[}21{]}{[}22{]}{[}24{]}. The analysis
establishes finite-time convergence guarantees and sample complexity
bounds for moving IoT service composition.

\textbf{Assumptions}: We assume the MDP satisfies standard regularity
conditions: (A1) state and action spaces are finite or compact, (A2)
policy parameterization is smooth with bounded gradients, (A3) value
function approximator uses linear or neural network function
approximation with bounded weights, (A4) step sizes satisfy
\(\sum \alpha_t = \infty\), \(\sum \alpha_t^2 < \infty\).

\textbf{Theorem 1 (Convergence Rate)}: Under assumptions A1-A4, the A2C
algorithm converges to an \(\epsilon\)-approximate stationary point with
sample complexity \(\mathcal{O}(\tilde{\epsilon}^{-2})\).

\emph{Proof Sketch}: Following {[}22{]}, we characterize error
propagation between actor and critic updates. The critic uses TD
learning with function approximation, introducing an approximation error
\(\varepsilon_{critic}\). The actor updates using the advantage
function, which introduces bias \(\varepsilon_{actor}\) from the value
function estimate. The total error bound combines these terms:

\[\| \nabla J(\theta) \| \leq \mathcal{O}(\varepsilon_{critic} + \sqrt{\varepsilon_{actor}} + \frac{1}{\sqrt{T}})\]

With linear function approximation in the critic and appropriate step
sizes, \(\varepsilon_{critic} = \mathcal{O}(1/\sqrt{T})\) and
\(\varepsilon_{actor} = \mathcal{O}(1/T)\), yielding the stated
complexity.

\textbf{Theorem 2 (Sample Complexity for Service Composition)}: For the
moving IoT service composition problem with state dimension \(d_s\) and
action dimension \(d_a\), the A2C algorithm achieves
\(\epsilon\)-optimal policy with sample complexity:

\[\mathcal{O}\left(\frac{d_s d_a}{\epsilon^2} \log\frac{1}{\delta}\right)\]

where \(\delta\) is the confidence parameter.

\emph{Proof Sketch}: The state space includes relative polar
coordinates (\(n \times 3\) for \(n\) services: \(r, \cos\theta,
\sin\theta\)), consumer position (2), distance matrix (\(n \times n\)),
and temporal features. This yields \(d_s = \mathcal{O}(n^2)\). The
action space includes \(n\) service selection actions plus a dummy
action for no valid service, giving \(d_a = \mathcal{O}(n)\).

\textbf{Corollary 1 (Scalability)}: The sample complexity scales
polynomially with the number of services \(n\), specifically
\(\mathcal{O}(n^2)\) for both state and action dimensions. This matches
the complexity of the original Double DQN while providing faster
convergence as demonstrated experimentally.

\textbf{Corollary 2 (Convergence Time)}: The expected convergence time
in wall-clock terms is:
\[T_{conv} = \mathcal{O}\left(\frac{1}{\eta_\pi \epsilon^2} + \frac{1}{\eta_V \epsilon^2}\right)\]

where \(\eta_\pi\) and \(\eta_V\) are the actor and critic learning
rates. With \(\eta_\pi = 0.0003\) and \(\eta_V = 0.0007\) as used in our
experiments, convergence to \(\epsilon = 0.01\) requires approximately
500 episodes for A2C Separate and 350 episodes for A2C Shared.

\hypertarget{network-architectures}{%
	\subsubsection{3.4 Network Architectures}\label{network-architectures}}

We use a dense fully connected neural network to train our model, following the exact architecture from prior work [1]:

\textbf{Neural Network Architecture}: The inputs to the network are the parameters defining a state whereas the outputs are the actions which the agent can take. Three hidden layers are used with 512 neurons each. The rectified linear unit (ReLU) activation function is used in all hidden layers. We use dropout with probability 0.5 on all hidden layers to reduce overfitting. The Q-learning's discount factor $\gamma$ is set to 0.9, whereas the Neural Network's learning rate is 0.001.

\hypertarget{relative-polar-coordinate-representation}{%
	\paragraph{3.4.1 Shared Architecture}\label{shared-architecture}}

A common feature extraction backbone followed by separate actor and
critic heads:

\begin{figure*}[t]
	\centering
	\caption{Architecture of A2C with Shared Network}
	\label{shared-architecture-diagram}
	\small
	\begin{verbatim}
Input Layer (State Vector + Distance Matrix)
    |
    v
+---------------------------------------+
|   Shared Feature Encoder             |
|   - FC(256) -> ReLU                    |
|   - FC(128) -> ReLU                    |
+---------------------------------------+
    |
    +-----------------------+-----------------------+
    v                       v
+-----------------+   +-----------------+
|   Actor Head    |   |   Critic Head   |
|   FC(action_dim)|   |   FC(1)         |
|   -> Softmax    |   |                 |
+-----------------+   +-----------------+
     Policy pi          Value V(s)
	\end{verbatim}
\end{figure*}

The shared encoder extracts spatio-temporal features from the state
representation including service positions, device trajectory, temporal
context, and the distance matrix for STR calculation.

\hypertarget{separate-architecture}{%
	\paragraph{3.4.2 Separate Architecture}\label{separate-architecture}}

Independent networks enable specialized processing:

\begin{figure*}[t]
	\centering
	\caption{Architecture of A2C with Separate Networks}
	\label{separate-architecture-diagram}
	\small
	\begin{verbatim}
Input State -> Actor Network                      Input State -> Critic Network
     |                                                    |
     v                                                    v
+----------------------+                      +----------------------+
| Trajectory Encoder   |                      | State Encoder        |
| LSTM(128)            |                      | FC(256) -> ReLU       |
+----------------------+                      +----------------------+
     |                                                    v
+----------------------+                      +----------------------+
| Policy Head          |                      | Value Head           |
| FC(64) -> ReLU       |                      | FC(64) -> ReLU       |
| FC(action_dim)      |                      | FC(1)                |
| -> Softmax           |                      +----------------------+
+----------------------+
     Policy pi               Value V(s)
	\end{verbatim}
\end{figure*}

The actor incorporates LSTM for trajectory-aware policy learning,
capturing temporal dependencies in service movement patterns. The critic
uses standard feedforward processing for value estimation.

\hypertarget{relative-polar-coordinate-representation}{%
	\subsection{3.5 Relative Polar Coordinate Representation}\label{relative-polar-coordinate-representation}}

Instead of velocity-based trajectory prediction, we transform raw GPS
coordinates into a relative polar coordinate representation that is
directly fed to the DRL agent. This transformation provides several
advantages: it encodes spatial relationships explicitly, is robust to
coordinate system variations, and provides a compact state
representation.

The polar coordinate transformation converts the absolute positions
into relative features:

\begin{equation*}
r_i = \sqrt{(x_i - x_c)^2 + (y_i - y_c)^2}, \qquad
\theta_i = \operatorname{atan2}(y_i - y_c, x_i - x_c)
\end{equation*}

The state vector for each service consists of three features:
\([r_i, \cos\theta_i, \sin\theta_i]\). This representation ensures that:
1. Distance \(r_i\) captures the proximity to the consumer
2. The angular components \(\cos\theta, \sin\theta\) capture directional
information while avoiding the discontinuity at \(\theta = \pi\)
3. The representation is translation-invariant with respect to the
consumer's position

\hypertarget{spatio-temporal-constraint-handling-with-str-based-capacity}{%
	\subsubsection{3.6 Spatio-Temporal Constraint Handling with STR-Based
		Capacity}\label{spatio-temporal-constraint-handling-with-str-based-capacity}}

Spatio-temporal constraints integrate through the reward function and
state representation using the hierarchical STR → Capacity model:

\textbf{Spatio-Temporal Validity (Discovery Zone)}: A service \(i\) is
valid (discoverable) at time \(t\) if it satisfies:
\[\text{Valid}(i, j, t) = \mathbb{1}(d_{ij}(t) \leq R_{comm})\]

where \(d_{ij}(t)\) is the Euclidean distance between service \(i\) and
consumer \(j\) at time \(t\), and \(R_{comm}\) is the communication
range (discovery zone). This is the fundamental spatio-temporal
constraint: the service must be physically located within the consumer's
wireless coverage area at the exact time of composition.

\textbf{STR-Based Service Quality}: Once validity is established, the
STR model ranks valid services:
\[STR(d_{ij}(t)) = \begin{cases} 1 & \text{if } d_{ij}(t) \leq R_c \\ e^{-k \cdot (d_{ij}(t) - R_c)} & \text{if } d_{ij}(t) > R_c \end{cases}\]

\textbf{Capacity Constraint (STR-Based)}: Composed services must meet
capacity requirements:
\[C_{ij}(t) = B \cdot \log_2(1 + STR(d_{ij}(t)) \cdot SNR_{max}) \geq C_{required}(t)\]

The reward function penalizes selecting invalid services (outside
discovery zone) with -1 reward. The STR-derived capacity serves as the
primary QoS metric following distance → STR → capacity hierarchy.

\begin{center}\rule{0.5\linewidth}{0.5pt}\end{center}

\hypertarget{experimental-setup}{%
